\section{Finite carriers for a stationary current}
\label{car:section}

We continue with the restricted HNN splitting, the free group $P$, and the
nonzero atomless stationary current
\[
 i_*\nu=j_*\nu=\mu
\]
provided by Proposition~\ref{stat:stationary}.  Let $T_P$ be the Cayley tree
of the fixed basis of $P$, let $R$ be its basis rose, and write
$\Omega=\operatorname{supp}\mu$.  Propositions~\ref{stat:transfers}
and~\ref{stat:windows} give a conull set $\mathcal U\subset\Omega$ of
lines with irrational endpoints, together with their unique aligned
bilateral trajectories through the canonical edge $e$ of the
Bass--Serre tree $\mathcal T$.  Transfers along finite parts of these
trajectories preserve measure.

Our objective is to identify the finitely generated noncyclic subgroups
that can fix an entire bilateral trajectory.  The window stabilizers
have bounded rank but can have arbitrarily large core graphs.  We keep
track of that distinction: after suppressing vertices of degree two,
the graph sizes are bounded, and the edges whose labels remain bounded
produce a finite family of subgroup carriers.  Maximal carriers return
to themselves under transfer, giving automorphisms to which we can
apply Brinkmann's theorem.

\subsection{Window cores and their limits}

We use the immersed-core description of finitely generated free
subgroups and the corresponding pullback description of intersections;
see Stallings \cite[Section~3.3, Theorem~5.5, and Section~7]{Stallings1983}.
The rank bounds and limiting carrier construction needed here are
proved below.

For a finitely generated nontrivial subgroup $K\le P$, put
\[
 C(K)=K\backslash\operatorname{Hull}_{T_P}(\Lambda K).
\]
This is its finite unbased core graph, with its label-preserving
immersion into $R$.  It has no vertices of degree one.  A circuit in an
immersed graph will mean a locally injective map from a subdivided
circle; it need not be an embedded simple cycle.

For each $n\ge1$, choose one representative of every $P$-orbit of
positively traversed windows
\[
 W=(E_{-n},\ldots,E_n),\qquad E_0=e,
\]
whose full pointwise stabilizer $K_W=\operatorname{Stab}_P(W)$ is noncyclic.
Define the finite graph
\[
 \Gamma_n=\coprod_W C(K_W),
 \qquad \Gamma_n\longrightarrow R.
\]
The choice is finite by Lemma~\ref{ann:window-bound}.  Set
$r_0=\bar r(P)$.  Since $\mu$ is nonzero and atomless, $P$ has rank at
least two, and hence $r_0\ge1$.

A \emph{state} of a labeled core is one of its vertices. In each
noncyclic core, call a vertex \emph{natural} if its degree is at least
three. A \emph{natural edge} is a maximal nonbacktracking path whose interior vertices
have degree two. Suppressing those interiors gives the \emph{natural
graph}.

\begin{lemma}[Window cores]
\label{car:window-cores}
The following properties hold.
\begin{enumerate}
\item The sum of the reduced ranks of the components of $\Gamma_n$
is at most $r_0$.  After suppressing degree-two interiors, $\Gamma_n$
has at most $3r_0$ unoriented edges and $2r_0$ vertices.
\item Every line in $\Omega$ has a label-preserving lift to $\Gamma_n$.
\item A ray with irrational endpoint has at most one starting state
for a lift among all components of $\Gamma_n$.
\end{enumerate}
The same statements hold for any other fixed finite positive window
shape containing its specified central edge, with the same bounds.
\end{lemma}

\begin{proof}
The reduced-rank assertion is Lemma~\ref{ann:window-bound}.  In a
noncyclic connected core, suppressing degree-two interiors leaves a
graph all of whose vertices have degree at least three.  If $E,V$ are
its edge and vertex counts and $r$ its rank, then
\[
 2E\ge3V,\qquad E-V=r-1.
\]
Consequently $E\le3(r-1)$ and $V\le2(r-1)$.  Sum these inequalities
over the components.  This bounds natural edges, not their subdivided
lengths.

We describe precisely what a state of a core represents.  If the
representative window is $W$ and $K=K_W$, a vertex $Kx$ of $C(K)$
represents the actual window $x^{-1}W$ through $e$.  This description
does not depend on the chosen representative of the coset.  It is
injective: equality $x^{-1}W=y^{-1}W$ says
$yx^{-1}\in\operatorname{Stab}_P(W)=K$, or equivalently $Kx=Ky$.
States in different representative cores cannot represent the same
actual window, because that would identify their $P$-orbits.

A ray starting at $Kx$ lifts in $T_P$ starting at $x$ and remaining in
$\operatorname{Hull}(\Lambda K)$.  Its label, read from the identity,
therefore has endpoint in $\Lambda(x^{-1}Kx)$.  If its endpoint is
irrational, Proposition~\ref{stat:windows} forbids two distinct represented
windows carrying that endpoint.  This proves uniqueness of the
starting state.  Continuation from a state is unique because the core
map is an immersion.

For a line in $\mathcal U$, the endpoints lie in the limit set of its
window stabilizer by Proposition~\ref{stat:windows}; hence its labels lift
to $\Gamma_n$.  Liftability to a fixed finite graph is closed in line
space.  Indeed, choose a vertex on a limiting line.  On each bounded
centered segment, sufficiently close lines have the same labels.
There are only finitely many possible lift states for the chosen
vertex, so a subsequence uses one state for arbitrarily long centered
segments.  The immersion then gives a lift of the whole line from
that state.  Since $\mathcal U$ is dense in $\Omega$, every
$\Omega$-line lifts.  This is the only extension from the conull set
to the entire support used here; we have not yet asserted that all
support lines admit transfers.

The proof applies to any fixed window shape, using the corresponding
rank bound and trajectory uniqueness.
\end{proof}

We shall also use Lemma~\ref{stat:window-finiteness}: for a fixed window
shape and a fixed ball in $T_P$, only finitely many actual windows
with noncyclic stabilizer have their stabilizer hull meeting that
ball.  In the preceding state description this is particularly
transparent.  If $x^{-1}\operatorname{Hull}(\Lambda K)$ contains a
vertex $y$ of the ball, then $Kxy$ is a vertex of the finite core
$C(K)$.  The finite choices of $y$ and $Kxy$ determine $Kx$, hence
the actual window $x^{-1}W$.

\begin{definition}
\label{car:small-graph}
Pass to a subsequence $n\to\infty$ on which the natural graph of
$\Gamma_n$ has a fixed abstract type, and identify those natural
graphs with one fixed finite graph.  Pass to a further subsequence
on which each natural edge has either a fixed label or length tending
to infinity.  For every edge of the latter kind, arrange convergence
of all finite prefixes of its label in both orientations.  The
\emph{small graph} $\Gamma^{\mathrm{sm}}$ consists of the natural
vertices and the natural edges with fixed labels, subdivided according
to those labels.
\end{definition}

Here is the order of these choices.  There are only finitely many
natural graph types with at most $3r_0$ edges and $2r_0$ vertices.
For each of the finitely many edges of a chosen type, take a
subsequence on which its lengths are bounded or tend to infinity.
In the bounded case there are only finitely many possible reduced
labels, so fix one.  In the divergent case diagonalize over the
prefix length, simultaneously in both orientations and on every
such edge.  These choices leave an infinite subsequence with
$n\to\infty$.  No bound on the lengths of the divergent edges is
assumed.

Figure~\ref{car:small-figure} illustrates the distinction between a
fixed natural graph and its fixed small subgraph.

\suppressfloats[t]
\begin{figure}[htbp]
\centering
\begin{tikzpicture}[scale=.93, every node/.style={font=\small},
  bounded/.style={draw=blue!55!black,very thick},
  longedge/.style={draw=black,densely dashed,thick}]
\node at (1.3,1.45) {$\Gamma_n$ (schematic)};
\draw[bounded] (0,0) .. controls (-1,1.15) and (-1,-1.15) .. (0,0);
\draw[bounded] (2.6,0) .. controls (3.6,1.15) and (3.6,-1.15) .. (2.6,0);
\draw[longedge] (0,0) -- node[above] {$w_n$}
 node[below=3pt] {$|w_n|\to\infty$} (2.6,0);
\fill (0,0) circle(2pt); \fill (2.6,0) circle(2pt);
\node[left] at (-.8,0) {$u$}; \node[right] at (3.4,0) {$v$};
\draw[-{Stealth[length=2mm]},thick] (4.15,0) -- (5.3,0);
\node[align=center] at (4.9,.65) {keep bounded\\edges};
\node at (8.1,1.45) {$\Gamma^{\mathrm{sm}}$};
\draw[bounded] (6.8,0) .. controls (5.8,1.15) and (5.8,-1.15) .. (6.8,0);
\draw[bounded] (9.4,0) .. controls (10.4,1.15) and (10.4,-1.15) .. (9.4,0);
\fill (6.8,0) circle(2pt); \fill (9.4,0) circle(2pt);
\node[left] at (6,0) {$u$}; \node[right] at (10.2,0) {$v$};
\node at (1.3,-1.2) {fixed rank, growing labels};
\node at (8.1,-1.2) {fixed labeled subgraph};
\end{tikzpicture}
\caption{A fixed natural graph can have arbitrarily long edge labels.
Here the bounded loops survive and the divergent bridge is deleted,
while both natural vertices remain. The small components in this
schematic example are cyclic; a noncyclic small component is not
assumed to exist.}
\label{car:small-figure}
\end{figure}

Each connected component $D$ of $\Gamma^{\mathrm{sm}}$, with a
chosen vertex, gives the subgroup of $P$ consisting of labels of
its based loops.  The map to $R$ is an immersion, so this subgroup
is isomorphic to $\pi_1D$.  Changing the vertex changes its
conjugacy class only.  Components of rank zero or one are allowed.
A bi-infinite reduced path entirely in $D$ lies in its unbased core,
and its labels have endpoints in a translate of the corresponding
subgroup limit set.

\subsection{A finite cofinal family of subgroups}

Let $\mathcal C$ be the set of $P$-conjugacy classes of finitely
generated noncyclic subgroups $Q\le P$ fixing a positively traversed
bilateral path through $e$ pointwise.  Such a path is unique: otherwise its
two finite windows would eventually diverge, contrary to
Proposition~\ref{stat:windows}.  Define
\[
 [Q]\preceq[M]
 \quad\Longleftrightarrow\quad
 Q\le aMa^{-1}\text{ for some }a\in P.
\]
We first verify that this really is a partial order, rather than
merely a preorder.

\begin{lemma}
\label{car:conjugate-containment}
A nontrivial finitely generated subgroup of a finite-rank free group
cannot be properly contained in a conjugate of itself.  Consequently
$\preceq$ is a partial order on $\mathcal C$.
\end{lemma}

\begin{proof}
Suppose $a^{-1}Ma\lneq M$.  No positive power $a^k$ belongs to $M$:
otherwise the nested inclusions
\[
 M\supseteq a^{-1}Ma\supseteq\cdots
 \supseteq a^{-k}Ma^k=M
\]
force equality at every stage.  Thus the Schreier vertices
$Ma^{-n}$, $n\ge0$, are distinct.  Fix a nontrivial reduced word
$w\in M$, of length $L$.  At each of those vertices the word $w$
labels a closed reduced path, since $a^{-n}wa^n\in M$.

The Schreier graph consists of its finite unbased core with trees
attached.  A closed reduced path of length $L$ has its starting
vertex at distance at most $L/2$ from the core: outside it, the
path must travel to the core and return along the same tree
attachment.  The $L/2$-neighborhood of the finite core has only
finitely many vertices.  It cannot contain all the distinct
$Ma^{-n}$, proving the first assertion.  The reverse direction of
conjugate containment follows by replacing $a$ by $a^{-1}$.

If $[Q]\preceq[M]\preceq[Q]$, suitable representatives give
\[
 Q\le aMa^{-1}\le abQ(ab)^{-1}.
\]
The inclusion between the first and last groups is equality by the
first assertion.  Both intermediate inclusions are therefore
equalities, and $[Q]=[M]$.  Reflexivity and transitivity are immediate.
\end{proof}

The following elementary graph bound supplies the compactness step.
It is important here that a core can have bridges.

\begin{lemma}[Bounded circuits in a core]
\label{car:circuits}
Let $C$ be a finite connected graph with no degree-one vertices and
with at least one edge.  Every edge of $C$ lies on an immersed circuit
of length at most $2|E(C)|$, where edges have length one and
$|E(C)|$ counts unoriented edges.
\end{lemma}

\begin{proof}
A nonbridge edge lies on a simple cycle of length at most $|E(C)|$.
If $a$ is a bridge, removing its interior leaves two components,
each containing a cycle.  Indeed, a component without cycles would
be a finite tree with at most one possible vertex of degree one,
which is impossible.  On each side choose a shortest path from
the endpoint of $a$ to a simple cycle, traverse the cycle, and
return along that path.  The resulting based path is reduced and
has length at most twice the number of edges on that side.
Combine these two paths with $a$ and $a^{-1}$.  The joins are
reduced, including the cyclic join, because the based paths remain
on their respective sides of $a$.  The resulting immersed circuit
contains $a$ and has length at most $2|E(C)|$.
\end{proof}

\begin{proposition}[Finite cofinal carriers]
\label{car:finite-carriers}
The noncyclic component subgroup classes of
$\Gamma^{\mathrm{sm}}$ belong to $\mathcal C$ and form a finite
cofinal subset: every $[Q]\in\mathcal C$ is contained, up to
$P$-conjugacy, in one of them.  In particular $\mathcal C$ has only
finitely many maximal elements, and they are exactly the maximal
members of this finite family.
\end{proposition}

\begin{proof}
Let $M$ be a noncyclic subgroup given by loops at a chosen vertex
of a small component.  These same loop labels occur at the
corresponding state of every selected $\Gamma_n$.  By the state
interpretation in Lemma~\ref{car:window-cores}, $M$ fixes an actual
centered window of radius $n$.  The windows obtained for different
$n$ have identical restrictions to any common shape, because
$M$ is noncyclic.  Their compatible restrictions form a bilateral
positive path fixed by $M$.  Thus $[M]\in\mathcal C$.

Now take $[Q]\in\mathcal C$, and let $C(Q)$ be its fixed unbased
core.  Put
\[
 D_Q=2|E(C(Q))|.
\]
For every selected $n$, the group $Q$ is contained in the full
stabilizer of its centered radius-$n$ window.  After translating to
one of our representatives, this inclusion gives a label-preserving
immersion
\[
 C(Q)\longrightarrow\Gamma_n.
\]
One way to see the map directly is to include the minimal
$Q$-invariant subtree in the minimal subtree of the containing
group and take quotients.  Conjugating the containing group merely
translates its subtree and does not change labels.

Choose $n$ so large that every divergent natural edge of
$\Gamma_n$ has length greater than $D_Q$.  Suppose the image of
some edge of $C(Q)$ meets the interior of one of those natural edges.  By
Lemma~\ref{car:circuits}, that source edge lies on an immersed
circuit of length at most $D_Q$.  Its image remains immersed.
At the degree-two interior of a natural edge there is no reduced
way to turn around, so the image circuit traverses that entire
natural edge.  This remains true if its chosen starting point is
in the interior: view the traversal cyclically.  Its length is
then greater than $D_Q$, a contradiction.

The whole image of $C(Q)$ therefore lies in
$\Gamma^{\mathrm{sm}}$.  Since it is connected, it lies in one
small component, with subgroup $M$.  The induced inclusion gives
$[Q]\preceq[M]$, and $M$ is noncyclic because it contains a
noncyclic subgroup.  Notice that no compatible choices of
conjugating elements were required as $n$ varied: once all long
edges exceed $D_Q$, the argument applies to every possible
label-preserving immersion.

There are finitely many small components.  A maximal member of
their noncyclic subgroup classes is maximal in all of
$\mathcal C$: any larger class lies below another member of the
finite family and hence would contradict its maximality there.
Conversely a maximal class in $\mathcal C$ equals a small class
above it.  This proves the final assertion.
\end{proof}

We have now replaced the possibly unbounded window cores by a
finite set of maximal subgroup carriers.  The next step uses
maximality to make a return transfer surjective.  If there are no
noncyclic small components, $\mathcal C$ is empty and all statements
about its maximal members below are vacuous.

\subsection{Return automorphisms}

Moving one edge forward along the path fixed by $Q$ and expressing
$Q$ in the new edge coordinates defines a map
\[
 F:\mathcal C\longrightarrow\mathcal C.
\]
This moves subgroup coordinates, not individual support lines;
it is defined for every class in $\mathcal C$.

\begin{proposition}[Atoroidal return]
\label{car:return}
The map $F$ is an order-preserving bijection with order-preserving
inverse.  It permutes the finite set of maximal classes of
$\mathcal C$.  For every representative $M$ of such a class, there
are an integer $p\ge1$ and an element $v\in G$ of height $p$ such
that
\[
 v^{-1}Mv=M.
\]
The automorphism $\alpha:M\to M$, $\alpha(m)=v^{-1}mv$, is
atoroidal: no nontrivial conjugacy class in $M$ is periodic under
$\alpha$.
\end{proposition}

\begin{proof}
Choose a frame $h$ for the next edge on the unique path fixed by
$Q$.  Transfer gives $h^{-1}Qh\le P$.  Another frame has the form
$hq$, $q\in P$, and conjugates this output inside $P$.
Conjugating $Q$ inside $P$ translates its path and gives the same
output class.  Thus $F$ is well defined.  Moving backward is its
inverse.  If $Q\le M$ after choosing representatives, both groups
fix the same path by noncyclic uniqueness; the same partial
conjugation preserves their inclusion.  This proves the order
assertions.

An order automorphism permutes maximal elements.  There are
finitely many of them by Proposition~\ref{car:finite-carriers},
so $F^p[M]=[M]$ for some $p\ge1$.  If their number is $N$, one
may choose $p\le N$.  Let $h$ be a frame for the $p$th forward
edge.  There exists $c\in P$ with
\[
 h^{-1}Mh=cMc^{-1}.
\]
Set $v=hc$.  Then $v$ has height $p$ and
$v^{-1}Mv=M$, which proves that the return map is an onto
automorphism.

Suppose $\alpha^k(w)=u^{-1}wu$ for some nontrivial $w\in M$,
$u\in M$, and $k\ge1$.  Then
\[
 g=v^ku^{-1}
\]
commutes with $w$.  Its height is $kp\ne0$, whereas $w$ has
height zero and infinite order.  A relation $w^a g^b=1$ forces
$b=0$ by height, then $a=0$.  Thus
$\langle w,g\rangle\cong\mathbb Z^2=BS(1,1)$, contrary to
BS-freeness.  This proves atoroidality without using hyperbolicity
of $G$ or the flaring statement under investigation.
\end{proof}

The element $v$ preserves the bilateral path fixed by $M$:
normalizing $M$ sends that path to another path fixed by $M$.
Both paths contain $ve$, so uniqueness in that edge's coordinates
identifies them.  It translates the path forward by $p$
edges.  In particular, the paths from $e$ to $v^be$ and to
$v^{-b}e$ are the concatenated positive and negative return
paths, respectively, and are reduced.

\subsection{Expansion and vanishing carrier mass}

For an element $w$ of a finitely generated free group with a chosen
basis, write $|w|$ for based word length and $\ell(w)$ for cyclic
length.  Brinkmann's theorem states that an automorphism $\alpha$
of a finite-rank free group with no nontrivial periodic conjugacy
class is hyperbolic: there exist $b_0\ge1$ and $\lambda_0>1$ such
that
\[
 \lambda_0|w|
 \le\max\{|\alpha^{b_0}(w)|,|\alpha^{-b_0}(w)|\}
 \quad\text{for every }w
\]
\cite[Theorem~1.1 and the definition of a hyperbolic
automorphism]{Brinkmann2000HyperbolicAutomorphisms}.  These are based lengths in a fixed
free basis.  We need both a stronger expansion factor and a
consequence for arbitrary currents; the latter uses a sum, not
a maximum.

\begin{lemma}[Expansion on elements and currents]
\label{car:expansion}
Let $M,v,\alpha$ be as in Proposition~\ref{car:return}, and choose
a free basis of $M$.  There exist $b\ge1$, an automorphism
$\beta=\alpha^b$, and $\lambda>4$ such that
\begin{align}
 \lambda|w|_M
 &\le\max\{|\beta(w)|_M,|\beta^{-1}(w)|_M\},
 &&w\in M,                                      \label{car:based}\\
 \lambda\ell_M(w)
 &\le\max\{\ell_M(\beta(w)),\ell_M(\beta^{-1}(w))\},
 &&w\in M.                                      \label{car:cyclic}
\end{align}
For every $M$-current $\sigma$ one has
\begin{equation}
 \lambda\|\sigma\|_M
 \le\|\beta_*\sigma\|_M+
       \|(\beta^{-1})_*\sigma\|_M.
 \label{car:current}
\end{equation}

The partial conjugations along the paths to $v^be$ and $v^{-b}e$
have domains
\[
 S_+=P\cap v^bPv^{-b},\qquad
 S_-=P\cap v^{-b}Pv^b,
\]
respectively.  These are finitely generated full window stabilizers,
and both contain $M$.  Their boundary transfers are defined on all
of $\partial^2\Lambda S_+$ and $\partial^2\Lambda S_-$.
The transfers $\Theta_+,\Theta_-$ restrict on $\partial^2\Lambda M$ to
$\beta,\beta^{-1}$, respectively, and preserve $\mu$ on their
irrational-endpoint domains.
\end{lemma}

\begin{proof}
Apply Brinkmann's theorem to $\alpha$, and put
$\gamma=\alpha^{b_0}$.  If
$|\gamma(w)|_M\ge\lambda_0|w|_M$, applying the inequality to
$\gamma(w)$ shows that its expanding direction must again be
forward: its backward image has length $|w|_M$, which is strictly
less than $\lambda_0|\gamma(w)|_M$ for $w\ne1$.  Induction gives
\[
 |\gamma^k(w)|_M\ge\lambda_0^k|w|_M.
\]
The analogous statement holds for a backward initial expanding
direction.  Choose $k\ge1$ with $\lambda_0^k>4$, and set
$b=b_0k$, $\beta=\alpha^b$, and $\lambda=\lambda_0^k$.
This proves~\eqref{car:based}; it also holds for $w=1$.
Apply it to $w^q$, divide by $q$, and let $q\to\infty$.
The identity $\lim_{q\to\infty}|w^q|_M/q=\ell_M(w)$ proves
\eqref{car:cyclic}.

For a counting current $\eta_w$, the cyclic inequality and
$\|\eta_w\|_M=\ell_M(w)$ imply~\eqref{car:current}, with the
maximum replaced by the sum.  The sum form is linear and therefore
holds for all finite nonnegative linear combinations of counting
currents.  We give the approximation argument extending it to
every current.

Fix $k\ge2$.  The directed graph with vertices the reduced words
of length $k-1$ and edges the reduced words of length $k$ has an
edge from a word's prefix to its suffix.  The masses
$\langle u,\sigma\rangle$ of directed segment cylinders based at the
identity and labeled by length-$k$ words $u$ form a
nonnegative circulation: both incoming and outgoing sums at a
vertex equal its length-$(k-1)$ cylinder mass.  Any such circulation
is a finite nonnegative sum of directed cycle circulations.
Indeed, follow positive-flow edges to find a directed cycle,
subtract the minimum edge flow on that cycle, and repeat; at
least one positive edge disappears at each step.

Write the circulation as $\sum_i a_i C_i$.  Each directed cycle
is the block circulation of a cyclically reduced word $w_i$,
including multiplicities if that word is a proper power.  Let
$C_i^{\mathrm{flip}}$ be the circulation obtained by reversing
words and inverting letters.  The counting current $\eta_{w_i}$
has circulation $C_i+C_i^{\mathrm{flip}}$.  Since $\sigma$ is
flip-invariant, the current
\[
 \sigma_k=\frac12\sum_i a_i\eta_{w_i}
\]
has exactly the same masses as $\sigma$ on all length-$k$
cylinders.  Summing over extensions gives equality on all shorter
word cylinders as well.  This also covers graph loops and
period-one words; proper powers have their prescribed counting
multiplicities.

As $k\to\infty$, the currents $\sigma_k$ converge to $\sigma$
on every cylinder, hence in the current topology.  Their norms
are equal to that of $\sigma$, since length-one masses already
agree.  Continuity of the norm and of the current maps induced
by $\beta^{\pm1}$ now proves~\eqref{car:current}.

Finally, the return paths are reduced as observed after
Proposition~\ref{car:return}.  Fixing their first and last edges
is equivalent to fixing the intervening segment, so their full
stabilizers in $P$ are exactly $S_+$ and $S_-$ above.  Since
$v$ normalizes $M$, both contain $M$.  Their partial conjugations
are $x\mapsto v^{-b}xv^b$ and $x\mapsto v^bxv^{-b}$, whose
restrictions are $\beta$ and $\beta^{-1}$.  The boundary and
measure assertions follow from Proposition~\ref{stat:windows}.
\end{proof}

\begin{proposition}[Zero mass on a maximal carrier]
\label{car:zero-mass}
For every maximal carrier $M$ supplied by
Proposition~\ref{car:finite-carriers},
\[
 \mu(\partial^2\Lambda M)=0.
\]
The same holds for every $P$-translate of this line set.
\end{proposition}

\begin{proof}
Use the boundary embedding of $M$ in $P$ to pull back the
restriction of $\mu$ to the relatively closed subset
$\partial^2\Lambda M$ of $\partial^2P$.  This gives an
$M$-current $\sigma$: invariance and flip-invariance come from
$\mu$, and the restriction is Radon because the boundary
embedding sends compact subsets of $\partial^2M$ to compact
subsets of $\partial^2P$.

Both endpoints are irrational outside a $\mu$-null set by
Proposition~\ref{stat:transfers}.  On those lines the return transfer
$\Theta_+$ is measure-preserving, maps
$\partial^2\Lambda M$ onto itself, and agrees with $\beta$.
Thus $\beta_*\sigma=\sigma$; the inverse has the same property.
Equation~\eqref{car:current} gives
\[
 \lambda\|\sigma\|_M\le2\|\sigma\|_M.
\]
Since $\lambda>4$, this forces $\sigma=0$.  Hence the restricted
measure is zero.  Translation invariance of $\mu$ gives the
last assertion.
\end{proof}

This is a measure conclusion, not a support conclusion: a closed
set of measure zero can still contain lines of $\Omega$.  The
next section uses the return expansion to exclude precisely those
possible support lines.
